\chapter{Программная реализация}\label{ch:impl}

\section{Структура библиотеки}

Библиотека реализована на языке C++ стандарта C++23 и состоит из 25 заголовочных файлов,
сгруппированных по подсистемам:

\begin{enumerate}
    \item \texttt{core} --- матричный тип и генератор случайных чисел с явным сидом;
    \item \texttt{autograd} --- механизм автоматического дифференцирования: узел графа,
          обратный проход и четырнадцать операций с правилами их дифференцирования;
    \item \texttt{data} --- загрузка выборки, стандартизация, стратифицированное
          разделение, унитарное кодирование меток и статистический анализ признаков;
    \item \texttt{nn} --- обучаемые блоки: полносвязный слой, функции активации, функции потерь, схемы
          инициализации весов и модель целиком;
    \item \texttt{optim} --- интерфейс оптимизатора и градиентный спуск с импульсом;
    \item \texttt{train} --- цикл обучения, история метрик по эпохам и метрики качества:
          доля правильных ответов, матрица ошибок, точность;
    \item \texttt{io} --- сохранение и загрузка модели в текстовом формате, включающем
          архитектуру сети, веса и параметры стандартизации;
    \item \texttt{viz} --- построение кривых обучения.
\end{enumerate}

Разбор аргументов командной строки вынесен в отдельный файл \texttt{apps/main.cpp}.
Листинги ключевых фрагментов приведены в приложении~\ref{app:listings}.

\section{Используемые библиотеки}

Для вспомогательных задач, не относящихся к предмету работы, применяются готовые
библиотеки: Eigen~3.4.0 для операций линейной алгебры, Matplot++~1.2.2 для построения
графиков, fmt~11.0.2 для форматирования вывода, rapidcsv~9.07 для чтения и записи файлов
формата CSV, CLI11~2.7.2 для разбора командной строки и Catch2~3.9.1 для модульного
тестирования.

\section{Интерфейс}

Программа оформлена в виде единого исполняемого файла с четырьмя подкомандами:
\texttt{describe} выполняет статистический анализ набора данных, \texttt{split} разделяет
его на обучающую и проверочную выборки, \texttt{train} обучает сеть и сохраняет модель,
\texttt{predict} загружает сохранённую модель, выполняет прогноз и оценивает его по
бинарной кросс-энтропии.

Архитектура и параметры обучения задаются аргументами командной строки: число и ширина
скрытых слоёв, функция активации, схема инициализации, функция потерь, число эпох, размер
порции данных, скорость обучения, коэффициент импульса, терпение ранней остановки и зерно
генератора случайных чисел.
